\documentclass[11pt]{article}
\usepackage{amsmath,amssymb}
\usepackage[margin=1in]{geometry}
\usepackage{hyperref}

\title{Ideas Parked --- Not Yet Merged}
\author{Staging document; each entry marks its eventual destination}
\date{}

\begin{document}
\maketitle

\noindent This document holds ideas discussed against Pattern Arithmetic
v15\_19, Inference, and Assembly, none of which have been written into
any of those three files yet. Each entry states its status and its
intended destination. Nothing here is settled; the point of this file
is to let that be true without leaving Paper 1 mid-edit.

\section{Destined for Pattern Arithmetic (untouched, remains open)}

\paragraph{Mix across differing directions.} No progress made. The
existing paragraph already states the gap precisely: Mix is defined
only among states sharing a direction $\hat n$, and what it should mean
across differing directions is unsettled, with a provisional
vector-sum recipe already tried and withdrawn. A concrete motivating
case surfaced in discussion: a single witnessed episode carrying
humidity, pressure, and air-speed readings, each naturally sitting at
its own direction on the sphere, needing to be read together. This
sharpens the existing open item with a real example; it does not
resolve it. \textbf{[Open, unchanged.]}

\section{Destined for Pattern Arithmetic (advances an existing open item)}

\paragraph{Location as a sort --- a third candidate.} The existing
paragraph names two candidates and adopts neither. A third was tested
here: dedicate separate, disjoint dimensions to content, to time, and
to location, rather than letting any of them share an axis.

\textbf{[Verified.]} Two shifts acting through disjoint content
dimensions commute exactly (checked numerically, not approximately:
time-then-location and location-then-time landed on the identical
point to full floating-point precision). Two shifts forced to share a
single content dimension do not commute, confirming the disjointness
is what matters, not rotation itself.

\textbf{[Verified.]} Content stays recoverable after both shifts:
projecting back onto the content dimensions alone reproduced the
original direction closely, while the time-axis and location-axis
each carried a clean, separately-readable amount of shift, with
nothing blended between them.

\textbf{[Open, unchanged.]} The encoding itself --- what specific
rotation a given time value or location corresponds to --- is not
stated anywhere. The geometry now supports it cleanly; the mapping
does not exist yet.

This is real, checkable content, not yet written as an addition to
the existing paragraph.

\section{Destined for Pattern Arithmetic (open item stays open, unchanged)}

\paragraph{Reducing a path by merging close actors.} The existing
paragraph asks for three things: a stated distance, a merge rule, and
a bound on how far a reduced path's endpoint may drift from the
unreduced one. Discussion here supplied a distance (see below) and
partially addressed the merge rule, but only for magnitude, not for
direction --- and direction is the harder half this item is actually
about.

\textbf{[Verified.]} Switching the strength calculation from the
companion paper's clip to Inference's already-established Hill form,
$r'=|z|/(K+|z|)$, changes nothing about the resulting direction
$\chi'=\arg(z)$. Checked directly on two constructed sequences with
identical content and opposite order: both gave the same angle under
either formula. The Hill substitution fixes evidence-count blindness
in strength; it does not touch the order-blindness this open item is
actually asking about.

\textbf{[Open, unchanged.]} No merge rule for direction exists. No
drift bound exists for either magnitude or direction. This item
remains exactly as open as it was before this discussion, now with
the magnitude half more precisely separated from the direction half.

\section{Destined for Inference (new mechanism)}

\paragraph{Candlestick summary: a richer merge rule for reference-atom
building.} Not a resolution of the Paper 1 item above --- a different
question. Paper 1's open item is about \emph{walking}: can a shortened
path still be walked to roughly the same endpoint. This is about
\emph{reading}: how to summarize a group of atoms for comparison
against a reference atom, which is Inference's own established
business (its reference-atom section), not $T$'s.

\textbf{[Verified.]} Open/close/high/low, computed the way a trading
candle is, distinguishes exactly the case flat Mix could not. Two
sequences, identical content, opposite order --- naive Mix gave the
identical angle ($47.5°$) for both. The candle's open and close swapped
between them ($20°\!\to\!160°$ versus $160°\!\to\!20°$), correctly
reading one as ending badly and the other as ending well. High and low
stayed identical for both, correctly capturing that only the
\emph{range} was shared, not the \emph{direction of travel}.

\textbf{[Settled, matches existing machinery.]} The scope rule ---
every atom belongs to exactly one time unit, no overlap between
buckets --- is not new; it is precisely the discrete-time-bucket
metadata already defined on a walk's own steps in the companion paper.

\textbf{[Open.]} This gives Inference a genuinely better reference-atom
merge rule than flat Mix. It is not proposed as a fix for Paper 1's
walking-specific open item, which remains untouched above. Whether
open/close/high/low itself needs a coherence check analogous to
$\bar R_c$ has not been considered.

\section{Destined for Inference (naming caution, not yet a mechanism)}

\paragraph{``Inference Mandelbrot'' is not Assembly's Mandelbrot, and
should not share its name.} Assembly's multi-zoom promotes an
unordered thing (a settled camp) and is fully handled by the Hill
bijection already verified there. Whatever would promote an
\emph{ordered} thing (a walked path, a movie, a time series) into one
higher-level atom is a different, harder problem --- it is the
direction half of the path-reduction item above, restated from a
different angle. \textbf{[Naming flag only.]} If this is ever built,
it wants its own name; reusing ``Mandelbrot'' for it would be the same
kind of collision already caught and corrected for ``collapse,''
``singularity,'' and the quantum swap test.

\section{Considered and found to need nothing filed}

For completeness: the reference-atom and precursor-witness mechanism,
the bundle-vector / uncapped $V_{\text{derived}}$ distinction from
Mix, and the $\pi=(x;w_1,\ldots,w_n)$ poem trajectory were all
discussed at length and all turned out to already be fully present in
Inference and Paper 1 respectively. Nothing from those threads is
parked here, since there was nothing left unresolved to park.

\end{document}
